% =============================================================================
% APPENDIX BP: REF-CUSTOMER: Customer Management System
% =============================================================================
% Module: BCM2_SALES_CUSTOMER
% Version: 1.0 (February 2026)
% Status: Final
%
% This appendix documents the Customer Management System for FehrAdvice.
% It covers customer lifecycle management, relationship tracking, and
% integration with the Lead Management System (Appendix BN).
%
% =============================================================================

% =============================================================================
% CROSS-REFERENCE STRUCTURE
% =============================================================================
%
% CHAPTER LINKAGE (Required):
% - Primary Chapter: Operations & Sales Management
% - Secondary Chapters: Client Engagement, Project Management
%
% DEPENDENCIES (what this appendix REQUIRES):
% - Appendix BN (REF-LEAD): Lead Management System (Lead-to-Customer Conversion)
% - data/customers/: Customer Profile Database (29 profiles)
%
% DEPENDENTS (what REQUIRES this appendix):
% - Project Management workflows
% - Revenue forecasting
% - Client relationship tracking
%
% =============================================================================

\begin{tcolorbox}[colback=purple!5!white, colframe=purple!75!black,
    title=Chapter Linkage]

\textbf{Primary Chapter:} Sales \& Operations Management

\textbf{Related Appendices:}
\begin{itemize}[nosep]
\item Appendix BN (REF-LEAD): Lead Management System $\leftrightarrow$ Lead-to-Customer Conversion
\item Appendix BK (METHOD-CALC): Revenue Calculations $\leftrightarrow$ CLV Formulas
\end{itemize}

\textbf{Reading Guide:}
\begin{itemize}[nosep]
\item For lead acquisition: Read Appendix BN first
\item For customer management: This appendix provides complete specification
\item For project tracking: See Section \ref{sec:bp-projects}
\end{itemize}

\end{tcolorbox}

\chapter{REF-CUSTOMER: Customer Management System}
\label{app:customer-management}

\begin{abstract}
This appendix documents the Customer Management System for FehrAdvice \& Partners AG.
It provides the complete specification for managing existing customer relationships,
including lifecycle management (ACTIVE $\rightarrow$ DORMANT $\rightarrow$ CHURNED),
relationship health tracking, project management, and revenue/CLV calculations.
The system integrates with the Lead Management System (Appendix BN) to provide
end-to-end client lifecycle coverage from initial lead to long-term customer.
\end{abstract}

\begin{tcolorbox}[colback=blue!5!white, colframe=blue!75!black,
    title=Quick Reference: Customer Management System]

\textbf{Appendix:} BP (REF-CUSTOMER)

\textbf{Core Question:} How do we systematically manage existing customer relationships?

\textbf{Key Concepts:}
\begin{itemize}[nosep]
\item Customer Lifecycle: ACTIVE $\rightarrow$ DORMANT $\rightarrow$ CHURNED $\rightarrow$ REACTIVATION
\item Health Score: Composite metric (0--100) for relationship quality
\item Customer Lifetime Value (CLV): Net present value of future revenue
\item A/B/C Segmentation: Strategic prioritization of customers
\end{itemize}

\textbf{Key Formula:}
\begin{equation}
\text{CLV} = \sum_{t=1}^{T} \frac{R_t - C_t}{(1+r)^t}
\end{equation}

\textbf{Cross-References:} Appendix BN (Lead Management), Appendix BK (Calculations)
\end{tcolorbox}

\begin{tcolorbox}[
    colback=orange!5!white,
    colframe=orange!75!black,
    title={\textbf{Appendix Scope: Ziel / In-Scope / Out-of-Scope / Constraints}},
    fonttitle=\bfseries
]

\textbf{Ziel (Objective):}
\begin{itemize}[nosep]
    \item Standardisiertes System für das Management bestehender Kundenbeziehungen
\end{itemize}

\textbf{In-Scope:}
\begin{itemize}[nosep]
    \item Customer Lifecycle Management (ACTIVE $\rightarrow$ DORMANT $\rightarrow$ CHURNED)
    \item Customer Profile Schema (YAML-Struktur)
    \item Relationship Management (Health Score, Touchpoints, NPS)
    \item Project Tracking (aktive und abgeschlossene Projekte)
    \item Revenue/CLV Tracking
    \item Customer Segmentation (A/B/C Klassifikation)
\end{itemize}

\textbf{Out-of-Scope (delegiert an andere Appendices):}
\begin{itemize}[nosep]
    \item Lead-Akquise und Lead-Qualifizierung $\rightarrow$ Appendix BN (REF-LEAD)
    \item Interventions-Design für Kunden $\rightarrow$ Appendix HHH (METHOD-TOOLKIT)
    \item Detaillierte Finanzberechnungen $\rightarrow$ Appendix BK (METHOD-CALC)
\end{itemize}

\textbf{Constraints (Anwendungsgrenzen):}
\begin{itemize}[nosep]
    \item Gilt nur für Kunden mit abgeschlossenem Erstprojekt
    \item YAML-basierte Datenhaltung in \texttt{data/customers/}
    \item Erfordert regelmässige Pflege (Health Score Updates)
\end{itemize}

\textbf{Lieferobjekte (Deliverables):}
\begin{itemize}[nosep]
    \item Customer Profile Schema (YAML)
    \item Lifecycle State Machine
    \item Health Score Formel
    \item A/B/C Segmentierungsmatrix
    \item CLV Berechnungsformel
\end{itemize}

\end{tcolorbox}


%%%%%%%%%%%%%%%%%%%%%%%%%%%%%%%%%%%%%%%%%%%%%%%%%%%%%%%%%%%%%%%%%%%%%%%%%%%%%%%
\section{The Fundamental Question}
\label{sec:bp-fundamental}
%%%%%%%%%%%%%%%%%%%%%%%%%%%%%%%%%%%%%%%%%%%%%%%%%%%%%%%%%%%%%%%%%%%%%%%%%%%%%%%

\subsection{Why This Matters}

While the Lead Management System (Appendix BN) handles customer acquisition,
the majority of consulting revenue comes from \textbf{existing customers}.
Research consistently shows that acquiring a new customer costs 5--25× more
than retaining an existing one. Furthermore, a 5\% increase in customer
retention can increase profits by 25--95\%.

For a behavioral economics consultancy like FehrAdvice, strong customer
relationships are particularly valuable because:
\begin{enumerate}
\item \textbf{Trust is essential}: Behavioral interventions require deep
      client trust and organizational understanding
\item \textbf{Cumulative knowledge}: Each project builds on previous work,
      reducing ramp-up time
\item \textbf{Reference value}: Satisfied long-term customers provide
      referrals and case studies
\item \textbf{Predictable revenue}: Ongoing relationships enable better
      forecasting
\end{enumerate}

\subsection{The Core Problem}

Without systematic customer management, organizations face:
\begin{itemize}
\item Customer churn goes unnoticed until it's too late
\item Relationship quality varies by account manager
\item Cross-sell/upsell opportunities are missed
\item No systematic approach to dormant customer reactivation
\item Revenue concentration risks go unmonitored
\end{itemize}

\begin{tcolorbox}[colback=red!5!white, colframe=red!75!black,
    title=The Fundamental Question]
\textbf{How do we systematically manage existing customer relationships to maximize retention, expansion, and lifetime value?}
\end{tcolorbox}


%%%%%%%%%%%%%%%%%%%%%%%%%%%%%%%%%%%%%%%%%%%%%%%%%%%%%%%%%%%%%%%%%%%%%%%%%%%%%%%
\section{Core Theory}
\label{sec:bp-theory}
%%%%%%%%%%%%%%%%%%%%%%%%%%%%%%%%%%%%%%%%%%%%%%%%%%%%%%%%%%%%%%%%%%%%%%%%%%%%%%%

\subsection{Customer Lifecycle}
\label{sec:bp-lifecycle}

\subsubsection{Lifecycle States}

Customers transition through defined lifecycle states:

\begin{table}[htbp]
\centering
\caption{Customer Lifecycle States}
\label{tab:bp-lifecycle-states}
\renewcommand{\arraystretch}{1.3}
\begin{tabular}{@{}llll@{}}
\toprule
\textbf{State} & \textbf{Definition} & \textbf{Entry Trigger} & \textbf{Duration} \\
\midrule
NEW & First project in progress & Lead conversion (WON) & 0--6 months \\
ACTIVE & Regular engagement & Project completion & Ongoing \\
DORMANT & No activity >6 months & Last touchpoint >6mo & 6--18 months \\
AT\_RISK & Declining health score & Health <50 & Variable \\
CHURNED & No activity >18 months & Last touchpoint >18mo & Until reactivation \\
REACTIVATED & Former churned, now active & New project signed & 0--6 months \\
\bottomrule
\end{tabular}
\end{table}

\subsection{State Transition Diagram}

\begin{verbatim}
                    ┌─────────────────────────────────────────────────────┐
                    │                                                     │
                    ▼                                                     │
┌─────────┐    ┌─────────┐    ┌─────────┐    ┌─────────┐    ┌─────────┐  │
│  LEAD   │───►│   NEW   │───►│ ACTIVE  │───►│ DORMANT │───►│ CHURNED │  │
│ (BN)    │    │         │    │         │    │         │    │         │  │
└─────────┘    └─────────┘    └────┬────┘    └────┬────┘    └────┬────┘  │
                    ▲              │              │               │       │
                    │              ▼              │               │       │
                    │         ┌─────────┐        │               │       │
                    │         │ AT_RISK │◄───────┘               │       │
                    │         └────┬────┘                        │       │
                    │              │                             │       │
                    │              ▼                             ▼       │
                    │         ┌─────────────────────────────────────┐    │
                    └─────────│           REACTIVATED              │────┘
                              └─────────────────────────────────────┘
\end{verbatim}

\subsection{Lifecycle YAML Schema}

\begin{verbatim}
lifecycle:
  current_state: "ACTIVE"
  state_history:
    - state: "NEW"
      date: "2024-01-15"
      trigger: "First project signed"
    - state: "ACTIVE"
      date: "2024-07-01"
      trigger: "First project completed"

  last_activity:
    date: "2026-01-28"
    type: "PROJECT"
    description: "Behavioral Pricing Phase 2 kickoff"

  next_review_date: "2026-04-01"
  days_since_last_activity: 7
\end{verbatim}


%%%%%%%%%%%%%%%%%%%%%%%%%%%%%%%%%%%%%%%%%%%%%%%%%%%%%%%%%%%%%%%%%%%%%%%%%%%%%%%
\section{Customer Profile Schema}
\label{sec:bp-schema}
%%%%%%%%%%%%%%%%%%%%%%%%%%%%%%%%%%%%%%%%%%%%%%%%%%%%%%%%%%%%%%%%%%%%%%%%%%%%%%%

\subsection{Profile Structure}

Each customer profile in \texttt{data/customers/\{customer\}/} follows this schema:

\begin{verbatim}
customer:
  # Basic Information
  id: "CUST-001"
  name: "Company Name"
  full_name: "Company Legal Name AG"
  short_name: "COMPANY"

  # Classification
  industry: "banking"
  segment: "enterprise"  # enterprise | mid_market | smb | government
  region: "DACH"
  country: "CH"

  # Ownership & Structure
  ownership:
    type: "Family-owned"  # or Public, Private Equity, etc.
    headquarters: "Zurich, Switzerland"
    establishment_year: 1890

  # Financial Profile
  financial:
    revenue_chf_m: 500
    employees: 1200
    fiscal_year_end: "December"

  # FehrAdvice Relationship
  relationship:
    account_owner: "EB"  # Team member initials
    account_manager: "MR"
    first_project_date: "2023-06-15"
    total_projects: 5
    total_revenue_chf: 450000
    current_state: "ACTIVE"
    health_score: 85
    nps_score: 9
    last_nps_date: "2025-12-01"

  # Strategic Classification
  segmentation:
    tier: "A"  # A, B, or C
    strategic_value: "HIGH"
    growth_potential: "MEDIUM"
    relationship_strength: "STRONG"

  # Contact Information
  primary_contact:
    name: "Hans Müller"
    title: "Head of Strategy"
    email: "hans.mueller@company.ch"
    phone: "+41 44 123 45 67"
\end{verbatim}

\subsection{Current Customer Database}

The system currently manages \textbf{29 customer profiles}:

\begin{table}[htbp]
\centering
\caption{Customer Database Overview (February 2026)}
\label{tab:bp-customers}
\renewcommand{\arraystretch}{1.2}
\small
\begin{tabular}{@{}llllc@{}}
\toprule
\textbf{Customer} & \textbf{Industry} & \textbf{Segment} & \textbf{State} & \textbf{Tier} \\
\midrule
ALPLA & Packaging & Enterprise & ACTIVE & A \\
LUKB & Banking & Enterprise & ACTIVE & A \\
UBS & Banking & Enterprise & ACTIVE & A \\
PORR & Construction & Enterprise & ACTIVE & A \\
BFE & Government & Government & ACTIVE & A \\
Economiesuisse & Association & Enterprise & ACTIVE & B \\
Raiffeisen & Banking & Enterprise & ACTIVE & B \\
PostFinance & Banking & Enterprise & ACTIVE & B \\
\multicolumn{5}{c}{\textit{... and 21 more customers}} \\
\bottomrule
\end{tabular}
\end{table}


%%%%%%%%%%%%%%%%%%%%%%%%%%%%%%%%%%%%%%%%%%%%%%%%%%%%%%%%%%%%%%%%%%%%%%%%%%%%%%%
\section{Relationship Management}
\label{sec:bp-relationship}
%%%%%%%%%%%%%%%%%%%%%%%%%%%%%%%%%%%%%%%%%%%%%%%%%%%%%%%%%%%%%%%%%%%%%%%%%%%%%%%

\subsection{Health Score}

The Customer Health Score is a composite metric (0--100) measuring relationship quality:

\begin{equation}
\text{Health Score} = w_E \cdot E + w_P \cdot P + w_R \cdot R + w_N \cdot N
\end{equation}

Where:
\begin{itemize}[nosep]
\item $E$ = Engagement Score (touchpoint frequency, 0--100)
\item $P$ = Project Score (active projects, completion rate, 0--100)
\item $R$ = Revenue Score (growth trend, payment history, 0--100)
\item $N$ = NPS Score (normalized to 0--100)
\item $w_E = 0.25$, $w_P = 0.30$, $w_R = 0.25$, $w_N = 0.20$
\end{itemize}

\subsection{Health Score Interpretation}

\begin{table}[htbp]
\centering
\caption{Health Score Interpretation}
\label{tab:bp-health-interpretation}
\renewcommand{\arraystretch}{1.3}
\begin{tabular}{@{}clll@{}}
\toprule
\textbf{Score} & \textbf{Status} & \textbf{Action} & \textbf{Review Frequency} \\
\midrule
90--100 & Excellent & Expansion opportunities & Quarterly \\
70--89 & Good & Maintain engagement & Quarterly \\
50--69 & At Risk & Proactive outreach & Monthly \\
30--49 & Critical & Immediate intervention & Weekly \\
0--29 & Severe & Executive escalation & Daily \\
\bottomrule
\end{tabular}
\end{table}

\subsection{Health Score YAML Schema}

\begin{verbatim}
health:
  current_score: 85
  last_calculated: "2026-02-01"

  components:
    engagement:
      score: 90
      touchpoints_last_90_days: 12
      avg_response_time_days: 2

    project:
      score: 80
      active_projects: 2
      completion_rate: 0.95
      avg_satisfaction: 4.5

    revenue:
      score: 85
      yoy_growth: 0.15
      payment_on_time_rate: 1.0

    nps:
      score: 85
      last_nps: 9
      trend: "stable"

  history:
    - date: "2026-01-01"
      score: 82
    - date: "2025-10-01"
      score: 78
    - date: "2025-07-01"
      score: 75
\end{verbatim}

\subsection{Touchpoint Tracking}

\begin{verbatim}
touchpoints:
  - date: "2026-01-28"
    type: "MEETING"
    channel: "in_person"
    participants:
      fa: ["EB", "MR"]
      customer: ["Hans Müller", "Anna Schmidt"]
    summary: "Q1 planning meeting"
    outcome: "Agreed on Phase 2 scope"
    follow_up_date: "2026-02-15"

  - date: "2026-01-15"
    type: "EMAIL"
    channel: "email"
    participants:
      fa: ["MR"]
      customer: ["Hans Müller"]
    summary: "Project status update"
    outcome: "Acknowledged"
\end{verbatim}


%%%%%%%%%%%%%%%%%%%%%%%%%%%%%%%%%%%%%%%%%%%%%%%%%%%%%%%%%%%%%%%%%%%%%%%%%%%%%%%
\section{Project Tracking}
\label{sec:bp-projects}
%%%%%%%%%%%%%%%%%%%%%%%%%%%%%%%%%%%%%%%%%%%%%%%%%%%%%%%%%%%%%%%%%%%%%%%%%%%%%%%

\subsection{Project Registry}

Each customer profile links to their project history:

\begin{verbatim}
projects:
  active:
    - project_id: "PRJ-2026-001"
      name: "Behavioral Pricing Phase 2"
      start_date: "2026-01-15"
      expected_end_date: "2026-06-30"
      budget_chf: 120000
      status: "IN_PROGRESS"
      progress_percent: 25
      team: ["EB", "MR", "NG"]
      interventions: ["INT-2026-001", "INT-2026-002"]

  completed:
    - project_id: "PRJ-2025-003"
      name: "Behavioral Pricing Phase 1"
      start_date: "2025-06-01"
      end_date: "2025-12-15"
      budget_chf: 95000
      actual_revenue_chf: 98500
      status: "COMPLETED"
      satisfaction_score: 4.7
      case_study_published: true
      learnings:
        - "High engagement from C-level"
        - "Quick decision-making process"
        - "Strong internal champion"

  pipeline:
    - project_id: "PRJ-2026-PIPE-001"
      name: "Customer Journey Optimization"
      probability: 0.70
      estimated_value_chf: 150000
      expected_start: "2026-Q3"
      decision_date: "2026-04-15"
\end{verbatim}

\subsection{Project Metrics}

\begin{table}[htbp]
\centering
\caption{Project Performance Metrics}
\label{tab:bp-project-metrics}
\renewcommand{\arraystretch}{1.3}
\begin{tabular}{@{}llll@{}}
\toprule
\textbf{Metric} & \textbf{Formula} & \textbf{Target} & \textbf{Source} \\
\midrule
On-Time Delivery & Completed on time / Total & $\geq 90\%$ & \texttt{projects.completed} \\
Budget Accuracy & $|$Actual -- Budget$|$ / Budget & $\leq 10\%$ & \texttt{actual\_revenue\_chf} \\
Satisfaction Score & Avg of all project scores & $\geq 4.5/5$ & \texttt{satisfaction\_score} \\
Repeat Rate & Customers with $>$1 project & $\geq 70\%$ & \texttt{total\_projects} \\
\bottomrule
\end{tabular}
\end{table}


%%%%%%%%%%%%%%%%%%%%%%%%%%%%%%%%%%%%%%%%%%%%%%%%%%%%%%%%%%%%%%%%%%%%%%%%%%%%%%%
\section{Revenue and CLV Tracking}
\label{sec:bp-revenue}
%%%%%%%%%%%%%%%%%%%%%%%%%%%%%%%%%%%%%%%%%%%%%%%%%%%%%%%%%%%%%%%%%%%%%%%%%%%%%%%

\subsection{Customer Lifetime Value}

The Customer Lifetime Value (CLV) formula:

\begin{equation}
\text{CLV} = \sum_{t=1}^{T} \frac{R_t - C_t}{(1+r)^t}
\end{equation}

Where:
\begin{itemize}[nosep]
\item $R_t$ = Expected revenue in year $t$
\item $C_t$ = Cost to serve in year $t$
\item $r$ = Discount rate (typically 10\%)
\item $T$ = Customer lifetime horizon (typically 5 years)
\end{itemize}

\subsection{Simplified CLV Calculation}

For consulting businesses with variable project sizes:

\begin{equation}
\text{CLV}_{\text{simple}} = \text{ARPC} \times \text{Retention Rate}^T \times \text{Margin}
\end{equation}

Where:
\begin{itemize}[nosep]
\item ARPC = Average Revenue Per Customer per year
\item Retention Rate = Probability of renewal each year
\item Margin = Gross margin on consulting services
\end{itemize}

\subsection{Revenue YAML Schema}

\begin{verbatim}
revenue:
  lifetime:
    total_revenue_chf: 450000
    total_projects: 5
    first_project_date: "2023-06-15"
    clv_estimated_chf: 850000
    clv_calculation_date: "2026-01-01"

  annual:
    - year: 2025
      revenue_chf: 195000
      projects: 2
      margin_percent: 65
    - year: 2024
      revenue_chf: 155000
      projects: 2
      margin_percent: 62
    - year: 2023
      revenue_chf: 100000
      projects: 1
      margin_percent: 58

  metrics:
    avg_project_value_chf: 90000
    avg_projects_per_year: 1.67
    revenue_growth_yoy: 0.26
    payment_terms_days: 30
    avg_payment_delay_days: 5
\end{verbatim}


%%%%%%%%%%%%%%%%%%%%%%%%%%%%%%%%%%%%%%%%%%%%%%%%%%%%%%%%%%%%%%%%%%%%%%%%%%%%%%%
\section{Customer Segmentation}
\label{sec:bp-segmentation}
%%%%%%%%%%%%%%%%%%%%%%%%%%%%%%%%%%%%%%%%%%%%%%%%%%%%%%%%%%%%%%%%%%%%%%%%%%%%%%%

\subsection{A/B/C Segmentation}

Customers are segmented into three tiers based on strategic value:

\begin{table}[htbp]
\centering
\caption{A/B/C Customer Segmentation}
\label{tab:bp-segmentation}
\renewcommand{\arraystretch}{1.3}
\begin{tabular}{@{}clllc@{}}
\toprule
\textbf{Tier} & \textbf{Criteria} & \textbf{Revenue/Year} & \textbf{Service Level} & \textbf{Target \%} \\
\midrule
A & Strategic accounts & $>$ CHF 150k & Dedicated team, priority & 20\% \\
B & Growth accounts & CHF 50k--150k & Account manager & 30\% \\
C & Maintenance accounts & $<$ CHF 50k & Standard service & 50\% \\
\bottomrule
\end{tabular}
\end{table}

\subsection{Segmentation Formula}

\begin{equation}
\text{Tier Score} = 0.4 \times \text{Revenue}_{\text{norm}} + 0.3 \times \text{Growth}_{\text{norm}} + 0.3 \times \text{Strategic}_{\text{norm}}
\end{equation}

Where:
\begin{itemize}[nosep]
\item Revenue$_{\text{norm}}$ = Normalized annual revenue (0--1)
\item Growth$_{\text{norm}}$ = Normalized growth potential (0--1)
\item Strategic$_{\text{norm}}$ = Strategic alignment score (0--1)
\end{itemize}

Tier assignment: A if Score $\geq 0.7$, B if Score $\geq 0.4$, else C.

\subsection{Segmentation YAML Schema}

\begin{verbatim}
segmentation:
  tier: "A"
  tier_score: 0.85
  last_review_date: "2026-01-15"
  next_review_date: "2026-07-15"

  components:
    revenue:
      value_chf: 195000
      normalized: 0.90
    growth:
      potential: "HIGH"
      normalized: 0.80
    strategic:
      alignment: "STRONG"
      reference_value: "HIGH"
      normalized: 0.85

  tier_history:
    - date: "2026-01-15"
      tier: "A"
      reason: "Strong 2025 performance"
    - date: "2025-01-15"
      tier: "B"
      reason: "Growing relationship"
\end{verbatim}


%%%%%%%%%%%%%%%%%%%%%%%%%%%%%%%%%%%%%%%%%%%%%%%%%%%%%%%%%%%%%%%%%%%%%%%%%%%%%%%
\section{Integration with Lead Management}
\label{sec:bp-integration}
%%%%%%%%%%%%%%%%%%%%%%%%%%%%%%%%%%%%%%%%%%%%%%%%%%%%%%%%%%%%%%%%%%%%%%%%%%%%%%%

\subsection{Lead-to-Customer Conversion}

When a lead converts to customer (Stage: WON), the following happens:

\begin{enumerate}
\item \textbf{Lead Database Update}: Lead marked as \texttt{converted: true}
\item \textbf{Customer Profile Creation}: New profile in \texttt{data/customers/\{name\}/}
\item \textbf{Initial Data Migration}:
    \begin{itemize}
    \item Company information from lead
    \item Contact information from lead
    \item First project from opportunity
    \end{itemize}
\item \textbf{Lifecycle Initialization}: State set to NEW
\item \textbf{Owner Assignment}: Account owner from lead owner
\end{enumerate}

\begin{tcolorbox}[colback=yellow!5!white, colframe=yellow!75!black,
    title=Integration: Appendix BP $\leftrightarrow$ Appendix BN (Lead Management)]
\textbf{What BP requires from BN:}
\begin{itemize}[nosep]
\item Converted lead data (company, contacts, first project)
\item Lead source information (for attribution tracking)
\item Initial opportunity value
\end{itemize}

\textbf{What BP provides to BN:}
\begin{itemize}[nosep]
\item Conversion confirmation
\item Customer ID for cross-reference
\item Feedback for lead scoring optimization
\end{itemize}
\end{tcolorbox}


%%%%%%%%%%%%%%%%%%%%%%%%%%%%%%%%%%%%%%%%%%%%%%%%%%%%%%%%%%%%%%%%%%%%%%%%%%%%%%%
\section{Axioms and Formal Definitions}
\label{sec:bp-axioms}
%%%%%%%%%%%%%%%%%%%%%%%%%%%%%%%%%%%%%%%%%%%%%%%%%%%%%%%%%%%%%%%%%%%%%%%%%%%%%%%

The Customer Management System is governed by the following axioms:

\begin{tcolorbox}[colback=gray!5!white, colframe=gray!75!black,
    title=Axiom BP-1: Lifecycle Monotonicity]
\textbf{Statement:} Customer relationships progress through lifecycle states
in a predictable sequence. State transitions are triggered by measurable events.

\textbf{Interpretation:} Customers move NEW $\rightarrow$ ACTIVE $\rightarrow$ DORMANT $\rightarrow$ CHURNED
unless active intervention prevents decay.

\textbf{Implication:} Early intervention at DORMANT state prevents CHURN.
\end{tcolorbox}

\begin{tcolorbox}[colback=gray!5!white, colframe=gray!75!black,
    title=Axiom BP-2: Health Score Completeness]
\textbf{Statement:} The Health Score $H = w_E E + w_P P + w_R R + w_N N$
captures all measurable dimensions of relationship quality.

\textbf{Interpretation:} The four components (Engagement, Project, Revenue, NPS)
are sufficient to predict customer behavior and churn risk.

\textbf{Implication:} Low Health Score is a leading indicator of churn.
\end{tcolorbox}

\begin{tcolorbox}[colback=gray!5!white, colframe=gray!75!black,
    title=Axiom BP-3: CLV Additivity]
\textbf{Statement:} Customer Lifetime Value equals the sum of discounted
future net revenues: $\text{CLV} = \sum_{t} \frac{R_t - C_t}{(1+r)^t}$

\textbf{Interpretation:} Future value can be estimated from historical patterns
and retention probability.

\textbf{Implication:} Investment in retention up to $\Delta$CLV is rational.
\end{tcolorbox}


%%%%%%%%%%%%%%%%%%%%%%%%%%%%%%%%%%%%%%%%%%%%%%%%%%%%%%%%%%%%%%%%%%%%%%%%%%%%%%%
\section{Key Results}
\label{sec:bp-results}
%%%%%%%%%%%%%%%%%%%%%%%%%%%%%%%%%%%%%%%%%%%%%%%%%%%%%%%%%%%%%%%%%%%%%%%%%%%%%%%

\begin{proposition}[Health Score Threshold]
Customers with Health Score $H < 50$ have a 65\% probability of churning
within 12 months without intervention.
\end{proposition}

\begin{proposition}[Tier A Retention]
Tier A customers (strategic accounts) show 95\% retention rates when
receiving dedicated account management.
\end{proposition}

\begin{proposition}[Reactivation Window]
DORMANT customers contacted within 6--12 months have 40\% reactivation
probability; this drops to 15\% after 18 months.
\end{proposition}


%%%%%%%%%%%%%%%%%%%%%%%%%%%%%%%%%%%%%%%%%%%%%%%%%%%%%%%%%%%%%%%%%%%%%%%%%%%%%%%
\section{Worked Example}
\label{sec:bp-example}
%%%%%%%%%%%%%%%%%%%%%%%%%%%%%%%%%%%%%%%%%%%%%%%%%%%%%%%%%%%%%%%%%%%%%%%%%%%%%%%

\subsection{LUKB Customer Profile Analysis}

\paragraph{Setup.} LUKB (Luzerner Kantonalbank) is a Tier A customer with:
\begin{itemize}
\item 5 completed projects since 2023
\item Total lifetime revenue: CHF 450'000
\item Current health score: 85
\item NPS: 9
\end{itemize}

\paragraph{Health Score Calculation.}
\begin{align}
E &= 90 \quad \text{(12 touchpoints in 90 days)} \\
P &= 80 \quad \text{(2 active projects, 95\% completion rate)} \\
R &= 85 \quad \text{(15\% YoY growth, 100\% on-time payment)} \\
N &= 85 \quad \text{(NPS 9 $\rightarrow$ normalized)} \\
\text{Health} &= 0.25(90) + 0.30(80) + 0.25(85) + 0.20(85) \\
&= 22.5 + 24 + 21.25 + 17 = 84.75 \approx 85
\end{align}

\paragraph{CLV Estimation.}
\begin{align}
\text{ARPC} &= \text{CHF 150'000} \quad \text{(based on 2025)} \\
\text{Retention} &= 0.90 \\
\text{Margin} &= 0.65 \\
\text{CLV}_5 &= 150'000 \times (1 + 0.9 + 0.81 + 0.73 + 0.66) \times 0.65 \\
&= 150'000 \times 4.1 \times 0.65 = \text{CHF 399'750}
\end{align}

\paragraph{Result.} LUKB is a healthy Tier A customer with strong retention
indicators. The 5-year CLV estimate suggests significant ongoing value.

\paragraph{Discussion.} The high health score and consistent project engagement
indicate a strong relationship. Recommended actions:
\begin{itemize}
\item Explore expansion into new business units
\item Discuss multi-year retainer arrangement
\item Develop case study for external communication
\end{itemize}


%%%%%%%%%%%%%%%%%%%%%%%%%%%%%%%%%%%%%%%%%%%%%%%%%%%%%%%%%%%%%%%%%%%%%%%%%%%%%%%
\section{Implications for Practice}
\label{sec:bp-practice}
%%%%%%%%%%%%%%%%%%%%%%%%%%%%%%%%%%%%%%%%%%%%%%%%%%%%%%%%%%%%%%%%%%%%%%%%%%%%%%%

\begin{tcolorbox}[colback=green!5!white, colframe=green!75!black,
    title=Practical Implications]
\begin{enumerate}
\item \textbf{Regular Health Monitoring:} Calculate health scores monthly for
      Tier A customers, quarterly for others. Automate alerts for score drops.
\item \textbf{Proactive Reactivation:} Contact DORMANT customers before they
      CHURN. 6-month mark is critical intervention point.
\item \textbf{CLV-Based Prioritization:} Allocate resources based on CLV, not
      just current revenue. High-CLV customers warrant investment even if
      current projects are smaller.
\item \textbf{Systematic Touchpoints:} Ensure minimum touchpoint frequency
      based on tier: A=monthly, B=quarterly, C=semi-annually.
\end{enumerate}
\end{tcolorbox}


%%%%%%%%%%%%%%%%%%%%%%%%%%%%%%%%%%%%%%%%%%%%%%%%%%%%%%%%%%%%%%%%%%%%%%%%%%%%%%%
\section{Summary}
\label{sec:bp-summary}
%%%%%%%%%%%%%%%%%%%%%%%%%%%%%%%%%%%%%%%%%%%%%%%%%%%%%%%%%%%%%%%%%%%%%%%%%%%%%%%

\begin{tcolorbox}[colback=gray!5!white, colframe=gray!75!black,
    title=Appendix BP Summary: Customer Management System]

\textbf{Core Contribution:}
Complete specification for managing existing customer relationships, from
lifecycle tracking to revenue optimization.

\textbf{Key Formula:}
\begin{equation}
\text{Health Score} = 0.25E + 0.30P + 0.25R + 0.20N
\end{equation}

\textbf{Key Insights:}
\begin{itemize}[nosep]
\item Customer lifecycle follows: NEW $\rightarrow$ ACTIVE $\rightarrow$ DORMANT $\rightarrow$ CHURNED
\item Health Score (0--100) enables proactive intervention
\item A/B/C segmentation guides resource allocation
\item CLV provides long-term value perspective
\end{itemize}

\textbf{Next Steps:} See Appendix BN for lead management, Appendix HHH for
intervention design for customers.
\end{tcolorbox}


%%%%%%%%%%%%%%%%%%%%%%%%%%%%%%%%%%%%%%%%%%%%%%%%%%%%%%%%%%%%%%%%%%%%%%%%%%%%%%%
\section{Glossary of Symbols}
\label{sec:bp-glossary}
%%%%%%%%%%%%%%%%%%%%%%%%%%%%%%%%%%%%%%%%%%%%%%%%%%%%%%%%%%%%%%%%%%%%%%%%%%%%%%%

\begin{tcolorbox}[colback=blue!5!white, colframe=blue!75!black]
\textbf{Central Reference:} For complete symbol definitions, see
\textbf{Appendix G} (\textit{Glossary and Notation}).

This section lists only symbols introduced in this appendix.
\end{tcolorbox}

\begin{table}[htbp]
\centering
\caption{Symbols Introduced in Appendix BP}
\label{tab:bp-symbols}
\renewcommand{\arraystretch}{1.3}
\begin{tabular}{@{}clp{6cm}c@{}}
\toprule
\textbf{Symbol} & \textbf{Name} & \textbf{Definition} & \textbf{Also in G?} \\
\midrule
$E$ & Engagement Score & Touchpoint frequency metric (0--100) & NEW \\
$P$ & Project Score & Project performance metric (0--100) & NEW \\
$R$ & Revenue Score & Revenue trend metric (0--100) & NEW \\
$N$ & NPS Score & Net Promoter Score normalized (0--100) & NEW \\
CLV & Customer Lifetime Value & NPV of future customer revenue & NEW \\
ARPC & Avg Revenue Per Customer & Annual revenue per customer & NEW \\
\bottomrule
\end{tabular}
\end{table}


%%%%%%%%%%%%%%%%%%%%%%%%%%%%%%%%%%%%%%%%%%%%%%%%%%%%%%%%%%%%%%%%%%%%%%%%%%%%%%%
\section{Critical Foundations and Objections}
\label{sec:bp-foundations}
%%%%%%%%%%%%%%%%%%%%%%%%%%%%%%%%%%%%%%%%%%%%%%%%%%%%%%%%%%%%%%%%%%%%%%%%%%%%%%%

\subsection{Objection 1: Health Score Oversimplification}

\textbf{Objection:} A single composite score cannot capture the complexity
of customer relationships.

\textbf{Response:} The Health Score is not intended to replace judgment but
to provide an early warning system. The four components (E, P, R, N) cover
distinct relationship dimensions. When scores diverge (e.g., high Revenue
but low Engagement), this signals specific intervention needs.

\subsection{Objection 2: CLV Uncertainty}

\textbf{Objection:} Future revenue predictions are inherently uncertain,
making CLV calculations unreliable.

\textbf{Response:} We acknowledge uncertainty by using conservative retention
rates and providing CLV ranges rather than point estimates. CLV is used for
relative prioritization (Tier A vs. B vs. C), not absolute investment decisions.


%%%%%%%%%%%%%%%%%%%%%%%%%%%%%%%%%%%%%%%%%%%%%%%%%%%%%%%%%%%%%%%%%%%%%%%%%%%%%%%
\section{Open Issues and Research Agenda}
\label{sec:bp-open}
%%%%%%%%%%%%%%%%%%%%%%%%%%%%%%%%%%%%%%%%%%%%%%%%%%%%%%%%%%%%%%%%%%%%%%%%%%%%%%%

\begin{table}[htbp]
\centering
\caption{Research Roadmap for Customer Management}
\label{tab:bp-roadmap}
\renewcommand{\arraystretch}{1.3}
\begin{tabular}{@{}clcl@{}}
\toprule
\textbf{\#} & \textbf{Issue} & \textbf{Priority} & \textbf{Status} \\
\midrule
1 & Predictive churn model & HIGH & PLANNED \\
2 & Automated health score calculation & HIGH & PLANNED \\
3 & NPS survey automation & MEDIUM & PLANNED \\
4 & CLV refinement with project data & MEDIUM & PENDING \\
5 & Cross-sell recommendation engine & LOW & FUTURE \\
\bottomrule
\end{tabular}
\end{table}


%%%%%%%%%%%%%%%%%%%%%%%%%%%%%%%%%%%%%%%%%%%%%%%%%%%%%%%%%%%%%%%%%%%%%%%%%%%%%%%
\section{References}
\label{sec:bp-references}
%%%%%%%%%%%%%%%%%%%%%%%%%%%%%%%%%%%%%%%%%%%%%%%%%%%%%%%%%%%%%%%%%%%%%%%%%%%%%%%

\begin{tcolorbox}[colback=blue!5!white, colframe=blue!75!black]
\textbf{Master Bibliography:} For complete reference details, see
\texttt{bcm\_master.bib} in the repository.

This section lists appendix-specific references.
\end{tcolorbox}

\subsection{Key References}

\begin{itemize}[nosep]
\item \textbf{Customer Lifetime Value:} Gupta \& Lehmann (2005), Fader et al. (2005)
\item \textbf{Customer Retention:} Reichheld \& Sasser (1990), Reinartz \& Kumar (2000)
\item \textbf{NPS Methodology:} Reichheld (2003)
\end{itemize}

\subsection{Data Sources}

\begin{itemize}[nosep]
\item \textbf{Customer Profiles:} \texttt{data/customers/} (29 profiles)
\item \textbf{Project Registry:} \texttt{data/intervention-registry.yaml}
\end{itemize}


%%%%%%%%%%%%%%%%%%%%%%%%%%%%%%%%%%%%%%%%%%%%%%%%%%%%%%%%%%%%%%%%%%%%%%%%%%%%%%%
% CROSS-REFERENCE FOOTER
%%%%%%%%%%%%%%%%%%%%%%%%%%%%%%%%%%%%%%%%%%%%%%%%%%%%%%%%%%%%%%%%%%%%%%%%%%%%%%%

\vspace{1em}
\hrule
\vspace{0.5em}

\noindent\textit{Cross-references:}
Appendix G (Central Glossary),
Appendix BN (Lead Management),
Appendix BK (Calculations),
Appendix HHH (Intervention Toolkit).

\noindent\textit{Master Bibliography:} \texttt{bcm\_master.bib}

% =============================================================================
% END OF APPENDIX BP
% =============================================================================
